% Shared methodology: supervised attack pipeline + unsupervised clustering cross-check.
\subsection{Attack Methodology}
\label{sec:eval-method}

We evaluate each privacy objective with a primary supervised attack and a complementary unsupervised cross-check.
The \emph{supervised} method trains classifiers on labeled data to predict hidden labels from the observables.
Performance is measured by AUC, with a chance baseline of $0.5$.
The \emph{unsupervised} method clusters the observables without using labels.
We evaluate agreement between the resulting clusters and the true labels using the Adjusted Rand Index (ARI) and Normalized Mutual Information (NMI).

We execute the generated policies on the \emph{Polygon~Amoy} testnet.
For each policy, a provider account submits the policy commitment using \texttt{commitPolicyWithGasPad}.
We then generate the Groth16 proof locally, acting as the TPS (Section~\ref{subsec:trust-model}).
A separate TPS account submits the access transaction using \texttt{requestAccessBoundWithGasPad}.
In the prototype, provider-to-TPS and requester-to-TPS exchanges are implemented as in-process calls and generate no additional on-chain artifacts.
We extract the observables defined in Section~\ref{sec:eval-threat} from the transaction receipts, calldata, and event logs.

\subsubsection{Supervised Attacks}
We evaluate each objective through binary classification \emph{tasks}, in which the attacker predicts a hidden label from the observables.
PCH tasks predict DUO-term presence or ICD-10 chapter membership (Section~\ref{sec:eval-pch}).
PI tasks distinguish between executions of two policies (Section~\ref{sec:eval-pi}).
SPP tasks distinguish between two categories of a structural or parameter property (Section~\ref{sec:eval-spp}).

For every task, we train two classifiers: logistic regression (LR) and random forest (RF), representing linear and non-linear models, respectively.
Each task is evaluated over \textbf{30 random seeds}, using the same protocol across PCH, PI, and SPP.

For each seed, we first randomly select one execution per policy to limit the influence of repeatedly executed policies on the AUC and its permutation test.
We then construct a \textbf{policy-disjoint $70/30$ train/test split}, ensuring that no policy appears in both sets.
Test AUC therefore measures attack performance on policies not seen during training.
Because policies are kept intact, exact per-class split proportions cannot always be achieved.
For each seed, we generate $512$ candidate splits and discard those that lack either class in the training or test set.
Among the valid candidates, we select the split that best approximates a $30\%$ test fraction within each class.

For each seed and classifier, we assess test AUC using a label-permutation test.
After training, we keep the test-set prediction scores fixed and randomly permute the test labels $R=500$ times.
We recompute AUC for each permutation to obtain its null distribution under no association between scores and labels.
The empirical $p$-value is
\begin{equation}
	p = \frac{k+1}{R+1},
	\label{eq:perm-p}
\end{equation}
where $k$ counts the permutations whose AUC is at least as large as the observed AUC.
The $+1$ correction prevents a zero $p$-value, with a minimum attainable value of $1/(R+1)$.
The result for a given seed and classifier is considered significant when $p<0.05$.

All tasks are binary, and the chance baseline for AUC is $0.5$ regardless of class balance.
Class imbalance may affect estimation uncertainty but does not change this baseline.

For each task, let $\mathrm{AUC}_s$ denote the test AUC for seed $s$, with $S=30$ seeds.
The results report the mean AUC and the AUC gap $G_{\mathrm{AUC}}$, defined as the mean absolute deviation from $0.5$:
\begin{equation}
	\mathrm{AUC}_{\text{mean}} = \frac{1}{S}\sum_{s=1}^{S}\mathrm{AUC}_s, \qquad
	G_{\mathrm{AUC}} = \frac{1}{S}\sum_{s=1}^{S}\left|\mathrm{AUC}_s-0.5\right|.
	\label{eq:auc-gap}
\end{equation}
Taking absolute deviations before averaging prevents values above and below $0.5$ from canceling.

The PASS/FAIL decision uses the mean AUC and the significant-seed fraction.
The latter is the fraction of seeds with $p<0.05$, reported as \%sig.
We classify each task as follows:
%
\begin{quote}
	\textbf{FAIL} if $\mathrm{AUC}_{\text{mean}}>0.55$ and at least $15$ of the $30$ seeds yield $p<0.05$.
	\textbf{PASS} otherwise.
\end{quote}
%
We refer to this rule as the \emph{leakage criterion}.
A PASS means that the evaluated attacks did not meet this criterion for the task.

The two conditions define an operational criterion based on effect size and consistency across splits.
With large test sets, even a small AUC increase above chance may be statistically significant.
The AUC threshold of $0.55$ prevents such small effects from triggering a FAIL solely on the basis of statistical significance.

The stability condition requires statistically significant results in at least half of the $30$ splits.
Because the splits overlap across seeds, the significant-seed fraction measures consistency across splits and is not interpreted as a combined $p$-value.

Because these thresholds are operational choices, we reapply the criterion over AUC thresholds of $0.52$--$0.60$ and significant-seed fractions of $0.3$--$0.7$.
At fractions of at least $0.5$, at most the two tasks discussed in the PCH and SPP result sections fail; lowering the fraction to $0.3$--$0.4$ adds at most two SPP contrasts, whereas raising it to $0.6$ yields none. No PI task fails under any of these threshold pairs.

In the evaluated run, the stability condition determines which tasks are flagged.
The AUC condition alone would flag four PCH and four SPP tasks.
In each objective, three of these tasks are significant in fewer than half of the splits.
The stability condition alone flags the same single task in each of PCH and SPP as the full leakage criterion.

In ten null reruns with policy-independent gas features, four contain a spurious PCH failure (a $3\%$ per-task frequency), while none produces an SPP failure.
We therefore examine the two observed failures in the corresponding result sections and keep them as FAIL.

The effectiveness of gas padding is assessed separately by the paired gas-padding ablation (Section~\ref{sec:eval-ablation}).
Using the same policies, features, splits, and classifiers, the attacks recover signal from the unpadded transcript, and that signal collapses after padding.
Thus, the leakage criterion summarizes residual attack performance, whereas the paired reduction provides the primary empirical evidence that padding suppresses the modeled gas channel.

The criterion is applied independently to LR and RF; a task fails if either classifier crosses both thresholds, and the result tables report the classifier with the higher mean AUC.

The classifiers receive all nine features stated in Section~\ref{sec:eval-threat}, including the seven that are constant by design.

\subsubsection{Unsupervised Clustering Attacks}
We run K-Means (with $K$ set to the number of classes, or $K=2$ for policy pairs) on the on-chain gas features, and score the recovered clusters against the true labels with ARI and NMI.
Values near $0$ mean the clusters carry no information about the hidden label; values near $1$ mean full leakage.

PCH clusters all executions globally, and SPP is clustered the same way on its $1{,}125$-execution view (Section~\ref{sec:eval-dataset}), once per annotation axis of Section~\ref{sec:eval-spp} with $K$ equal to that axis's bucket count.
At these sample sizes, ARI and NMI are read directly, each averaged over $20$ K-Means seeds.
To test whether an ARI exceeds chance, the cluster assignments are kept fixed, the labels are shuffled $200$ times, and $p_{\mathrm{ARI}}$ follows Equation~\eqref{eq:perm-p} with $R=200$, so its smallest attainable value is $1/201 \approx 0.005$.

For PI, clustering is scored per pair, so each side of a pair must contribute enough executions; this requires two safeguards.
First, the PI clustering analysis uses a dedicated set of \emph{balanced policy pairs} constructed from depth-layer executions (Section~\ref{sec:eval-dataset}).
We evaluate $127$ pairs, each containing $30$ executions per policy.
The policy pairs used in the supervised PI attack (Section~\ref{sec:eval-pi}) are unsuitable for per-pair clustering.
Most pairs contain only one execution per policy, while some contain $1$ execution for one policy and $30$ for the other.
These small or highly imbalanced samples make clustering scores unreliable indicators of policy distinguishability.
The supervised attack is unaffected, because it pools the executions of each pair type.
Second, NMI is \emph{permutation-debiased}. Even after balancing, each side of a pair holds only $30$ executions. At that scale NMI is inflated even without signal, whereas ARI is not (Section~\ref{sec:eval-nmi-sensitivity}), so ARI is read directly and NMI must be debiased.
The cluster assignments are computed once (K-Means does not see the labels) and the observed NMI is averaged over $20$ K-Means seeds; the pair labels are then shuffled $200$ times and the NMI recomputed against the fixed assignments on each shuffle, giving the null distribution of NMI under zero true signal.
The debiased NMI is the observed NMI minus the mean of this null distribution, and the empirical $p$-value $p_{\mathrm{NMI}}$ follows Equation~\eqref{eq:perm-p} with $R=200$, where $k$ counts the shuffles whose NMI meets or exceeds the observed value.
The cross-check flags a pair only if its debiased NMI is positive \emph{and} significant ($p_{\mathrm{NMI}}<0.05$).
